\documentclass{situationist-preprint}

\title{The Scope of A Priori Knowledge}
\msid{EPI-12-a}
\msfamily{EPI-12}
\msdiscipline{Epistemology}
\mstype{Principal manuscript}
\mskeywords{a priori; rationalism; empiricism; constitutive rules; intellectual seemings; logic; arithmetic; analyticity; safety}

\begin{document}
\maketitle

\begin{abstract}
Is there knowledge independent of experience? Rationalists have claimed a priori knowledge of logic and mathematics, of the geometry of space, of the principle that every event has a cause, of the fundamental truths of metaphysics and of the first principles of morality; radical empiricists have claimed that there is none. We argue that there is a priori knowledge, and that it extends to logic, arithmetic and some conceptual truths, and no further. A priori justification has two sources: the competent exercise of rules constitutive of the concepts a proposition contains, and intellectual seemings. Both can justify; only the first yields knowledge. Using the analysis of knowledge as safe, justified, true belief, we show that a belief formed by derivation from no premises according to constitutive rules of non-defective concepts is safe, and that a belief resting on an intellectual seeming about a matter that the concepts involved do not settle is not. Logic and arithmetic, given that arithmetic is logic plus Hume's principle, and analytic truths are conceptually grounded; the geometry of physical space, the uniformity of nature, universal causation and substantive metaphysics are not. The rationalist's intuitions can justify, but they cannot make knowledge.
\end{abstract}

\section{The problem}

A belief is justified a priori if its justification does not rest on experience. Experience may be needed to acquire the concepts the belief involves, as one must have seen colors to think about them; what matters is that experience plays no evidential role \citep{Kitcher1980,Burge1993}. The question whether there is a priori knowledge has two parts: whether there is any, and how far it extends.

The rationalist tradition gave generous answers to both. Descartes and Leibniz held that the existence of God, the nature of mind and the principle of sufficient reason are knowable by reason alone. \citet{Kant1787} held that geometry, arithmetic and the principle that every alteration has a cause are synthetic and known a priori, and made the explanation of this the task of the \emph{Critique}. Twentieth-century rationalists retreated from Kant's list but kept the faculty: \citet{BonJour1998} argued that rational insight into necessary truths is the only way to avoid scepticism, and \citet{Bealer1998} that intuitions are evidence for the propositions of philosophy as observations are for those of science.

Empiricists have narrowed or denied the domain. The logical empiricists restricted the a priori to analytic truths, true in virtue of meaning, and held that geometry as a theory of physical space is empirical \citep{Reichenbach1958}. \citet{Quine1951} argued that no statement is immune from revision in the light of experience, and that the analytic--synthetic distinction is untenable; \citet{Putnam1962} drew from the history of geometry the lesson that even apparent necessities can be overturned; \citet{Devitt2005} concluded that there is no a priori knowledge. A third position, developed by \citet{Boghossian1996} and \citet{Peacocke1992}, returned to a restricted rationalism: some propositions are knowable a priori because accepting them, or inferring according to them, is constitutive of possessing the concepts they involve.

We defend a version of the third position, and we give an argument for its scope that neither side has supplied. The argument rests on the analysis of knowledge as safe, justified, true belief of EPI-01 and on the distinction between justification and knowledge of EPI-06. Rationalists are right that intellectual seemings justify; empiricists are right that, outside the conceptual domain, they do not yield knowledge.

\section{Framework}

\begin{definition}[A priori justification and knowledge]\label{d:apriori}
A subject's justification for believing $p$ is \emph{a priori} if no experience figures in it as evidence, that is, as a node or root of its support other than in making possible the possession of the concepts involved. Her knowledge that $p$ is \emph{a priori} if it is knowledge, in the sense of EPI-01, whose justification and basis are a priori.
\end{definition}

\begin{definition}[Constitutive rule]\label{d:const}
A rule of inference is \emph{constitutive} of a concept if a general, defeasible disposition to infer according to it is part of what it is to possess the concept.
\end{definition}

Modus ponens is constitutive of the conditional, conjunction introduction and elimination of conjunction, and the inference from \enquote{$x$ is a vixen} to \enquote{$x$ is a female fox} of the concept of a vixen \citep{Peacocke1992}. Definition~\ref{d:const} does not require that every possessor of the concept infer according to the rule in every case; it requires that the disposition be part of the competence that possessors share, and that deviant possessors be deviant relative to it.

\begin{definition}[Conceptual grounding]\label{d:ground}
A proposition is \emph{conceptually grounded} if it can be derived from no premises by rules constitutive of the concepts it contains, including the logical concepts.
\end{definition}

A concept is \emph{defective} if its constitutive rules, taken together, are not truth-preserving under any assignment of a semantic value to it. The standard example is Prior's \enquote{tonk}, whose introduction rule is that of disjunction and whose elimination rule is that of conjunction \citep{Prior1960}; the naive concept of set, constituted by unrestricted comprehension, is another \citep{Russell1903}.

\begin{principle}[Two sources]\label{p:sources}
A priori justification has two sources. A subject who derives $p$ from no premises by competent exercise of constitutive rules has a priori justification for believing $p$; and, by the principle of EPI-19, a subject to whom it intellectually seems that $p$, on reflection and independently of experience, has prima facie a priori justification for believing $p$.
\end{principle}

\ref{p:sources} supplies what EPI-05 left open. The tree structure defended there concerns empirical justification, in which logical and arithmetical principles act mostly as rules of inference, edges of the tree rather than nodes. When such principles are themselves believed, their justification is either an inference from no premises, which the first source covers, or a root in an intellectual seeming, which the second covers. Both sources are internalist in the sense of EPI-06: whether one has competently derived $p$, and whether $p$ seems true to one on reflection, depend on one's mental states.

\begin{assumption}[Constitution]\label{a:const}
The constitutive rules of a non-defective concept are truth-preserving: the concept's semantic value is whatever makes them so.
\end{assumption}

Assumption~\ref{a:const} is the core of the conceptual-role account of the a priori \citep{Boghossian1996,Boghossian2000}. Possessing a concept is being disposed to use it by certain rules; the concept refers to whatever, if anything, makes those rules valid. If nothing does, the concept is defective, and the possessor's dispositions, though they fix what she means, do not make her inferences good.

\begin{assumption}[Indifference of non-conceptual seemings]\label{a:indiff}
If $p$ is not conceptually grounded, whether $p$ seems true to a subject on reflection is fixed by her concepts and cognitive constitution, which would be the same if $p$ were false; and the seemings of this kind include many that are false.
\end{assumption}

The first clause is a consequence of Definition~\ref{d:apriori}: an a priori seeming is not produced by experience of the matter it concerns, so if the concepts do not settle whether $p$, nothing in the production of the seeming does. The second clause is a matter of record. It seemed self-evident to the best minds that the geometry of space is Euclidean, that every set of things satisfying a condition forms a set, that simultaneity is absolute and that every event has a sufficient cause. Each seeming was firm, reflective and independent of experience, and each is false or, in the last case, contradicted by our best physics.

\begin{principle}[Safety for necessary truths]\label{p:safe}
A belief is safe only if, in close cases, the propositions the subject believes on the same basis are true.
\end{principle}

\ref{p:safe} is the amendment to the safety condition that EPI-01 accepts for necessary truths, following \citet{Pritchard2005}: since a necessary truth is true in every world, the safety of a belief in it must be assessed over the similar propositions the same basis would deliver.

\section{The argument}

\begin{lemma}[Safety of derivation]\label{l:deriv}
A belief in a conceptually grounded proposition, formed by competent derivation from no premises according to constitutive rules of non-defective concepts, is safe.
\end{lemma}

\begin{proof}
Consider the propositions the subject would believe, in close cases, on the same basis: by competent derivation from no premises according to the constitutive rules of the same non-defective concepts. By Assumption~\ref{a:const}, those rules are truth-preserving, and a derivation from no premises by truth-preserving rules has a true conclusion. So all such propositions are true, and by \ref{p:safe} the belief is safe.
\end{proof}

\begin{lemma}[Unsafety of non-conceptual seemings]\label{l:seem}
A belief in a proposition that is not conceptually grounded, resting on an intellectual seeming alone, is not safe.
\end{lemma}

\begin{proof}
The basis is a seeming, on reflection and independently of experience, about a matter that the concepts involved do not settle. By Assumption~\ref{a:indiff}, seemings of this kind are produced in the same way whether their contents are true or false, and many of them are false. Close cases in which the subject, or a subject with her constitution, believes on the same basis include cases in which the propositions so believed are false: had she been a geometer in 1850, or a set theorist in 1900, the same kind of seeming would have produced a false belief. By \ref{p:safe}, the belief is not safe.
\end{proof}

\begin{theorem}[Scope]\label{t:scope}
There is a priori knowledge. A proposition can be known a priori if it is conceptually grounded by the rules of non-defective concepts, and only if it is conceptually grounded.
\end{theorem}

\begin{proof}
If $p$ is conceptually grounded by non-defective concepts, a subject who competently derives it from no premises has a priori justification by \ref{p:sources}, holds a safe belief by Lemma~\ref{l:deriv}, and believes a truth; by EPI-01 she knows $p$, and her knowledge is a priori. Since subjects do derive simple logical truths in this way, there is a priori knowledge. Conversely, suppose $p$ is not conceptually grounded. By \ref{p:sources}, the only a priori basis available for believing it is an intellectual seeming; by Lemma~\ref{l:seem}, a belief resting on it is not safe, and so is not knowledge.
\end{proof}

It remains to say which propositions are conceptually grounded.

\begin{proposition}[Logic]\label{p:logic}
The truths of classical first-order logic, and of second-order logic so far as they are derivable in its deductive system, are conceptually grounded.
\end{proposition}

\begin{proof}
Each is derivable from no premises by the introduction and elimination rules of the logical constants, which are constitutive of them. Classical logic is the correct logic (LOG-02), so these rules are those of non-defective concepts; second-order logic is logic (LOG-09), so its rules are constitutive of logical concepts.
\end{proof}

The restriction in the second clause is required by incompleteness. Second-order logic with its standard semantics has truths that no derivation reaches, and these are not conceptually grounded in the sense of Definition~\ref{d:ground}; whether they are known, and how, is not settled here.

\begin{proposition}[Arithmetic]\label{p:arith}
The theorems of second-order Peano arithmetic are conceptually grounded.
\end{proposition}

\begin{proof}
Hume's principle, that the number of $F$s is the number of $G$s if and only if the $F$s and the $G$s are equinumerous, is constitutive of the concept of cardinal number \citep{Wright1983,HaleWright2001}, and MTH-04 argues that it is logical in the relevant sense. The axioms of second-order Peano arithmetic are derivable from it in second-order logic (Frege's theorem). Hence its theorems are derivable from no premises by rules constitutive of the logical concepts and of the concept of number.
\end{proof}

\begin{proposition}[Conceptual truths]\label{p:concept}
Some conceptual truths outside logic and arithmetic, such as that vixens are foxes and that whatever is red is colored, are conceptually grounded; which truths are, is in some cases indeterminate.
\end{proposition}

\begin{proof}
The inferences from \enquote{vixen} to \enquote{female fox} and from \enquote{red} to \enquote{colored} are constitutive of the concepts, so the propositions are derivable from no premises by constitutive rules. The indeterminacy follows from LNG-04: the analytic--synthetic distinction is vague, and so, for some rules, it is vague whether the disposition to infer by them is part of possessing the concept.
\end{proof}

The qualifier \enquote{some} in the thesis has two sources. The first is the vagueness just noted. The second is semantic externalism. Many concepts, such as that of water, have their reference fixed in part by the environment (LNG-08), and their constitutive rules settle little about what falls under them; that water is H$_2$O is necessary but not conceptually grounded, and it is known a posteriori \citep{Kripke1980}. Kripke's examples of the contingent a priori raise questions of their own, which EPI-12-b takes up.

\begin{corollary}[What lies outside]\label{c:outside}
The following are not known a priori: the geometry of physical space; the uniformity of nature; the principle that every event has a cause; the principle of sufficient reason and other substantive metaphysical principles; moral principles beyond those that are conceptual; and claims of metaphysical possibility inferred from conceivability where the concepts involved do not reveal the essences of their referents.
\end{corollary}

\begin{proof}
None is conceptually grounded, so by Theorem~\ref{t:scope} none is known a priori. In detail: the geometry of physical space is an empirical question, and general relativity answers it against Euclid \citep{Reichenbach1958,Putnam1962}. The uniformity of nature cannot be justified a priori, since no rule of prediction is reliable in every world (EPI-03). That every event has a cause, and that every fact has a sufficient reason, are not derivable from the concepts of event, cause, fact and reason, whose constitutive rules are consistent with uncaused events. Moral facts, on MTE-01, are natural facts about idealized reactions, and moral intuition is, on MTE-08, a faculty comparable to perception, so moral knowledge beyond the conceptual is not a priori. Finally, where the concept of a thing does not reveal its essence, as MND-05 argues for phenomenal concepts, the conceivability of a scenario does not establish its possibility, and the inference cannot yield a priori knowledge of possibility.
\end{proof}

Corollary~\ref{c:outside} is the precise content of \enquote{less than rationalists have claimed}. The rationalist's intellectual seemings about these matters are not worthless: by \ref{p:sources} and EPI-19 they confer prima facie justification, and a metaphysician who believes the principle of sufficient reason because it seems true to her on reflection may be justified in believing it. She does not know it, even if it is true, since a basis of that kind would easily have delivered false beliefs.

\section{Objections}

\begin{objection}[Revisability]
Every belief, including the beliefs of logic, is revisable in the light of experience \citep{Quine1951}. Proposals to revise logic in response to quantum mechanics show that even logic is held hostage to empirical theory. A priori knowledge, if it is knowledge immune from empirical revision, does not exist.
\end{objection}

\begin{reply}
Definition~\ref{d:apriori} does not require immunity from revision. A priori justification is justification in which experience plays no evidential role, and it is defeasible like any other: experience may provide defeaters, as the discovery of a paradox did for naive set theory. What a priori knowledge cannot be is knowledge whose positive support is experiential. As to quantum logic, LOG-13 argues that quantum mechanics revises probability theory, not logic, and that the distributive law survives. The mere logical possibility of revising a logical law under empirical pressure shows that we are fallible, not that our logical knowledge rests on experience.
\end{reply}

\begin{objection}[No understanding--assent links]
\citet{Williamson2007} argued that for any purported constitutive rule there are competent speakers who reject it: logicians who reject modus ponens for certain conditionals, philosophers who deny that every vixen is a vixen because they hold that empty terms make sentences false. If competence does not require acceptance, there are no constitutive rules, and the conceptual-role account of the a priori fails.
\end{objection}

\begin{reply}
Definition~\ref{d:const} requires a general, defeasible disposition shared by the competent, not acceptance by every possessor in every case. The deviant logician possesses the conditional by sharing the competence of the community whose usage fixes its meaning, and by being disposed to infer according to modus ponens in the ordinary cases; his theoretical rejection of the rule in special cases is a revision of belief, not an absence of the concept. Williamson's cases show that the possession of a concept does not guarantee the acceptance of its constitutive rules, which we grant. They do not show that the rules do not fix the concept's reference, which is what Assumption~\ref{a:const} requires, nor that competent derivation by them is an unsafe basis, which is what Lemma~\ref{l:deriv} denies.
\end{reply}

\begin{objection}[Self-defeat]
Rationalists argue that any argument against the a priori must itself rely on a priori reasoning: its premises, or at least the inferential connections between them, must be known independently of experience \citep{BonJour1998}. If intellectual seemings do not yield knowledge, the present argument, whose premises are philosophical and rest on reflection, does not yield knowledge either.
\end{objection}

\begin{reply}
The inferential connections in the argument are logical, and logic is within the scope of the a priori by Proposition~\ref{p:logic}. The premises are of two kinds. Assumption~\ref{a:const} and \ref{p:safe} are claims about concepts and about knowledge which we take to be close to conceptual; to that extent they fall within the a priori. Assumption~\ref{a:indiff} is in part a historical claim, and its support is empirical. The argument does not claim to be known a priori as a whole, and it does not need to be; it needs to be known, or justified, which it can be on a mixed basis. The self-defeat charge applies only to an empiricism that denies all a priori knowledge, which we reject.
\end{reply}

\begin{objection}[Mathematics beyond arithmetic]
Set theory, analysis and geometry as a branch of pure mathematics are known, and not by experience. If arithmetic is the only mathematics within the scope of the a priori, mathematical knowledge is either empirical or not knowledge.
\end{objection}

\begin{reply}
Pure geometry, as the theory of the structures satisfying its axioms, is conditional and conceptually grounded in the way logic is: what follows from the axioms is derivable from no premises by conditional proof. Analysis, so far as it is developed from the theory of the natural numbers by definitions, inherits the status of arithmetic. Set theory beyond what is conceptually grounded in this sense is a harder case. Its stronger axioms are accepted in part for their fruitfulness and their consequences, a form of extrinsic justification that MTH-06 relies on in arguing that the continuum hypothesis is false given new, independently motivated axioms. Knowledge resting on extrinsic justification is not a priori in the sense of Definition~\ref{d:apriori} when the consequences that support an axiom include applications, and it is not purely conceptual in any case. That such knowledge exists is consistent with the thesis, which concerns the a priori and not mathematics as a whole. MTH-02 holds that abstract objects are known through reliable belief-forming methods; Lemma~\ref{l:deriv} identifies one such method and explains its reliability.
\end{reply}

\begin{objection}[The existence of numbers]
MTH-01 holds that the existence of numbers is established by the indispensability argument, which is empirical. If arithmetic is known a priori, the existence of numbers, which it entails, is known a priori, and the indispensability argument is idle.
\end{objection}

\begin{reply}
A proposition can be knowable in more than one way, and knowledge on one basis does not make another basis idle. If Hume's principle is constitutive of the concept of number and the concept is not defective, arithmetic and the existence of numbers are knowable a priori. Whether the concept of number is defective is not settled by its constitutive rules, any more than the defectiveness of the naive concept of set was evident to those who possessed it; the indispensability argument provides independent, empirical support for the claim that it is not. The two routes are complementary.
\end{reply}

\section{What the result does not establish}

Theorem~\ref{t:scope} depends on Assumptions~\ref{a:const} and~\ref{a:indiff}. If there were a source of a priori justification, beyond derivation and seemings, whose reliability could be explained for non-conceptual matters, the argument would not exclude it; we know of no candidate. The theorem does not settle which concepts are defective, a question that may need experience, as the history of set theory shows. It does not deny that intellectual seemings about non-conceptual matters justify, and so does not convict rationalist metaphysicians of irrationality; it denies that they know. Finally, the theorem does not establish that conceptually grounded knowledge is easy to acquire. A proposition may be derivable from no premises only by a long derivation, and a subject who reaches it by intuition rather than derivation does not have the safe basis that Lemma~\ref{l:deriv} describes.

\section{Conclusion}

There is a priori knowledge, and its scope is set by its source. A belief derived from no premises by the constitutive rules of non-defective concepts is safe, since those rules are truth-preserving by the nature of the concepts; logic, arithmetic and some conceptual truths are derivable in this way. A belief resting on an intellectual seeming about a matter the concepts do not settle is not safe, since such seemings arise whether or not they are true and have often been false. The rationalist's seemings about space, causation, sufficient reason and morality can justify; they cannot make knowledge.

\begin{thebibliography}{99}\small
\bibitem[Bealer(1998)]{Bealer1998} Bealer, G. (1998). Intuition and the autonomy of philosophy. In M.~DePaul and W.~Ramsey (Eds.), \emph{Rethinking Intuition}. Rowman \& Littlefield.
\bibitem[Boghossian(1996)]{Boghossian1996} Boghossian, P.~A. (1996). Analyticity reconsidered. \emph{No\^us}, 30(3), 360--391.
\bibitem[Boghossian(2000)]{Boghossian2000} Boghossian, P.~A. (2000). Knowledge of logic. In P.~Boghossian and C.~Peacocke (Eds.), \emph{New Essays on the A Priori}. Oxford University Press.
\bibitem[BonJour(1998)]{BonJour1998} BonJour, L. (1998). \emph{In Defense of Pure Reason}. Cambridge University Press.
\bibitem[Burge(1993)]{Burge1993} Burge, T. (1993). Content preservation. \emph{The Philosophical Review}, 102(4), 457--488.
\bibitem[Devitt(2005)]{Devitt2005} Devitt, M. (2005). There is no a priori. In M.~Steup and E.~Sosa (Eds.), \emph{Contemporary Debates in Epistemology}. Blackwell.
\bibitem[Hale and Wright(2001)]{HaleWright2001} Hale, B. and Wright, C. (2001). \emph{The Reason's Proper Study: Essays towards a Neo-Fregean Philosophy of Mathematics}. Oxford University Press.
\bibitem[Kant(1787)]{Kant1787} Kant, I. (1781/1787). \emph{Kritik der reinen Vernunft}.
\bibitem[Kitcher(1980)]{Kitcher1980} Kitcher, P. (1980). A priori knowledge. \emph{The Philosophical Review}, 89(1), 3--23.
\bibitem[Kripke(1980)]{Kripke1980} Kripke, S.~A. (1980). \emph{Naming and Necessity}. Harvard University Press.
\bibitem[Peacocke(1992)]{Peacocke1992} Peacocke, C. (1992). \emph{A Study of Concepts}. MIT Press.
\bibitem[Pritchard(2005)]{Pritchard2005} Pritchard, D. (2005). \emph{Epistemic Luck}. Oxford University Press.
\bibitem[Prior(1960)]{Prior1960} Prior, A.~N. (1960). The runabout inference-ticket. \emph{Analysis}, 21(2), 38--39.
\bibitem[Putnam(1962)]{Putnam1962} Putnam, H. (1962). It ain't necessarily so. \emph{The Journal of Philosophy}, 59(22), 658--671.
\bibitem[Quine(1951)]{Quine1951} Quine, W.~V.~O. (1951). Two dogmas of empiricism. \emph{The Philosophical Review}, 60(1), 20--43.
\bibitem[Reichenbach(1958)]{Reichenbach1958} Reichenbach, H. (1958). \emph{The Philosophy of Space and Time}. Dover. (Original work published 1928.)
\bibitem[Russell(1903)]{Russell1903} Russell, B. (1903). \emph{The Principles of Mathematics}. Cambridge University Press.
\bibitem[Williamson(2007)]{Williamson2007} Williamson, T. (2007). \emph{The Philosophy of Philosophy}. Blackwell.
\bibitem[Wright(1983)]{Wright1983} Wright, C. (1983). \emph{Frege's Conception of Numbers as Objects}. Aberdeen University Press.
\end{thebibliography}

\end{document}
